\uuid{oo1Q}
\exo7id{4053}
\titre{exo7 4053}
\auteur{quercia}
\organisation{exo7}
\datecreate{2010-03-11}
\isIndication{false}
\isCorrection{false}
\chapitre{Développement limité}
\sousChapitre{Equivalents}
\cpge{mpsi-mp2i.analyse-asymptotique.a, pcsi.analyse-asymptotique.a, ptsi.analyse-asymptotique.a, mpsi-mp2i.analyse-asymptotique.b, pcsi.analyse-asymptotique.b, ptsi.analyse-asymptotique.b, mpsi-mp2i.analyse-asymptotique.c, pcsi.analyse-asymptotique.c, ptsi.analyse-asymptotique.c, bcpst1.limites-continuite.b, bcpst1.developpements-limites.a, bcpst1.developpements-limites.b}
\module{Analyse}
\niveau{L1}
\difficulte{}

\contenu{
\begin{enumerate}
    \item \question{Soit $f(x) = \ln(1+x) - x$.
  }
    \item \question{\'Etudier $f$.}
    \item \question{Chercher un équivalent simple de $f$ en $0$.}
    \item \question{Soit $(x_n)$ une suite de réels telle que $f(x_n) =  o(1/n)$.
        Montrer que $nx_n^2 \to 0$ lorsque $n\to\infty$.}
\end{enumerate}
}
